\chapter{Communication Complexity — Information Terms}
%%% generated by the blueprint pipeline from TCSlib.CommunicationComplexity.NewmanTheorem.FuncDisjointnessLowerBound.InformationTerms
\section{Overview}

This module contains the information-theoretic half of the linear lower bound for set
disjointness (Rao--Yehudayoff, Thm.~6.13): the special-coordinate information terms
$I(A_T : S \mid T\,A_{<T}\,B_{\ge T}\,\mathcal{D})$ and
$I(B_T : S \mid T\,A_{\le T}\,B_{>T}\,\mathcal{D})$, the fact that protocol duality swaps
the two terms, and the chain-rule upper bound $\ell \log 2 / n$ on each term for a protocol
of communication $\ell$.  The declaration documented here is the shared engine behind the
duality statements: a transfer lemma for conditional mutual information under the
sample-space and protocol dualities.

%%%%%%%%%%%%%%%%%%%%%%%%%%%%%%%%%%%%%%%%%%%%%%%%%%%%%%%%%%%%%%%%%%%%%%%%%%%%%%%
\section{Declarations}

\begin{theorem}[Information transfer along protocol duality]
\lean{CommunicationComplexity.Functions.Disjointness.RandomizedLowerBound.condMutualInfo_message_dualProtocol_eq}
\label{CommunicationComplexity.Functions.Disjointness.RandomizedLowerBound.condMutualInfo_message_dualProtocol_eq}
\leanok
\uses{CommunicationComplexity.Functions.Disjointness.RandomizedLowerBound.HardSample, CommunicationComplexity.Functions.Disjointness.RandomizedLowerBound.ProtocolType, CommunicationComplexity.Functions.Disjointness.RandomizedLowerBound.disjointCondMeasure, CommunicationComplexity.Functions.Disjointness.RandomizedLowerBound.disjointCondMeasure_measurePreserving_dualHardSample, CommunicationComplexity.Functions.Disjointness.RandomizedLowerBound.dualConditioningValue, CommunicationComplexity.Functions.Disjointness.RandomizedLowerBound.dualConditioningValue_injective, CommunicationComplexity.Functions.Disjointness.RandomizedLowerBound.dualHardSample, CommunicationComplexity.Functions.Disjointness.RandomizedLowerBound.dualProtocol, CommunicationComplexity.Functions.Disjointness.RandomizedLowerBound.dualProtocolTranscriptMap, CommunicationComplexity.Functions.Disjointness.RandomizedLowerBound.dualProtocolTranscriptMap_injective, CommunicationComplexity.Functions.Disjointness.RandomizedLowerBound.identDistrib_self_comp_measurePreserving, CommunicationComplexity.Functions.Disjointness.RandomizedLowerBound.message, CommunicationComplexity.Functions.Disjointness.RandomizedLowerBound.message_dualProtocol_dualHardSample, ProbabilityTheory.IdentDistrib.condMutualInfo_eq_finite, ProbabilityTheory.condMutualInfo_comp_right_conditioning_of_injective}
Let $p$ be a protocol on the hard sample space, let $X, X' : \Omega \to \{0,1\}$ be two
Boolean variables and let $Z, Z' : \Omega \to [n] \times \{0,1\}^n \times \{0,1\}^n$ be two
conditioning variables on the hard sample $\Omega = \texttt{HardSample}\,n$.  Suppose that
$X'$ and $Z'$ evaluated on the dual sample are $X$ and the recoded $Z$ on the original
sample, i.e.\ $X'(\sigma(\omega)) = X(\omega)$ and
$Z'(\sigma(\omega)) = \delta(Z(\omega))$ for every $\omega$, where $\sigma$ is the
hard-sample duality \texttt{dualHardSample} and $\delta$ the conditioning recoding
\texttt{dualConditioningValue}.  Then, under the disjoint-conditioned measure $\mu_{\mathcal D}$,
the conditional mutual information between $X'$ and the transcript of the dual protocol
$\bar p$ given $Z'$ equals that between $X$ and the transcript of $p$ given $Z$:
\[
  I\bigl[X' : M_{\bar p} \mid Z' ;\ \mu_{\mathcal D}\bigr]
  \;=\;
  I\bigl[X : M_{p} \mid Z ;\ \mu_{\mathcal D}\bigr],
\]
where $M_p = \texttt{message}\,n\,p$ denotes the transcript of $p$.
\end{theorem}
